\section{Theoretical Analysis}
\label{sec:theory}

We formalize three foundational properties of \texttt{Hetero-KVCache-Optimizer}:
(1)~the bounded error of attention under Sink+Tail position preservation,
(2)~the strict $O(1)$ memory independence from sequence length, and
(3)~the bandwidth-limited latency model for edge deployment.

\subsection{RoPE Position Invariance Lemma}
\label{subsec:rope-lemma}

Modern LLMs encode positional information via Rotary Position Embeddings (RoPE)~
\cite{su2024roformer}.  For a token at absolute position $m$, RoPE transforms the
query and key vectors as
\begin{equation}
\mathbf{q}_m = \mathbf{R}_{\Theta,m} \, \mathbf{q}, \qquad
\mathbf{k}_n = \mathbf{R}_{\Theta,n} \, \mathbf{k},
\end{equation}
where $\mathbf{R}_{\Theta,m}$ is a block-diagonal rotation matrix parameterized by a
frequency basis $\Theta = \{\theta_i\}_{i=1}^{d/2}$.  The attention score between
positions $m$ and $n$ depends only on the relative distance $(m-n)$:
\begin{equation}
\mathbf{q}_m^{\top} \mathbf{k}_n
= \mathbf{q}^{\top} \mathbf{R}_{\Theta,m}^{\top} \mathbf{R}_{\Theta,n} \mathbf{k}
= \mathbf{q}^{\top} \mathbf{R}_{\Theta,n-m} \mathbf{k}.
\end{equation}

\begin{lemma}[RoPE Position Invariance under Sink+Tail]
\label{lem:rope}
Let $\mathcal{S}$ be the set of Sink tokens with absolute positions
$\{0, \dots, L_s-1\}$ and $\mathcal{T}$ be the set of Tail tokens with absolute
positions $\{N-L_t, \dots, N-1\}$.  Let $\mathcal{M} = [0,N-1] \setminus
(\mathcal{S} \cup \mathcal{T})$ denote the evicted middle tokens.  If the KV cache
retains the original absolute position IDs $m$ and $n$ for every token in
$\mathcal{S} \cup \mathcal{T}$, then for any query at position
$p \in \mathcal{S} \cup \mathcal{T}$ the attention output computed over the pruned
cache satisfies
\begin{equation}
\bigl\| \text{Attn}_{\text{pruned}}(p) - \text{Attn}_{\text{full}}(p) \bigr\|_2
\;\le\;
\frac{1}{\sqrt{d_k}} \sum_{n \in \mathcal{M}}
\bigl| e^{\mathbf{q}_p^{\top} \mathbf{k}_n / \sqrt{d_k}} \bigr|,
\label{eq:rope-bound}
\end{equation}
where $d_k$ is the head dimension.
\end{lemma}

\begin{proof}
The full attention output at position $p$ is
\begin{equation}
\text{Attn}_{\text{full}}(p)
= \sum_{n=0}^{N-1} \alpha_{p,n} \, \mathbf{v}_n,
\qquad
\alpha_{p,n}
= \frac{e^{\mathbf{q}_p^{\top} \mathbf{k}_n / \sqrt{d_k}}}
        {\sum_{j=0}^{N-1} e^{\mathbf{q}_p^{\top} \mathbf{k}_j / \sqrt{d_k}}}.
\end{equation}
The pruned attention output omits the middle tokens $\mathcal{M}$ and
renormalizes over $\mathcal{S} \cup \mathcal{T}$:
\begin{equation}
\text{Attn}_{\text{pruned}}(p)
= \sum_{n \in \mathcal{S} \cup \mathcal{T}}
  \tilde{\alpha}_{p,n} \, \mathbf{v}_n,
\qquad
\tilde{\alpha}_{p,n}
= \frac{e^{\mathbf{q}_p^{\top} \mathbf{k}_n / \sqrt{d_k}}}
        {\sum_{j \in \mathcal{S} \cup \mathcal{T}}
         e^{\mathbf{q}_p^{\top} \mathbf{k}_j / \sqrt{d_k}}}.
\end{equation}
Because the absolute position IDs are preserved, every retained token keeps its
exact rotation matrix $\mathbf{R}_{\Theta,n}$; therefore
$\mathbf{q}_p^{\top} \mathbf{k}_n$ is unchanged for all $n \in \mathcal{S} \cup \mathcal{T}$.
Let $Z = \sum_{j=0}^{N-1} e^{\mathbf{q}_p^{\top} \mathbf{k}_j / \sqrt{d_k}}$ and
$\tilde{Z} = \sum_{j \in \mathcal{S} \cup \mathcal{T}}
e^{\mathbf{q}_p^{\top} \mathbf{k}_j / \sqrt{d_k}}$.
The difference between the full and pruned outputs is
\begin{align}
&\text{Attn}_{\text{full}}(p) - \text{Attn}_{\text{pruned}}(p) \nonumber \\
&= \sum_{n \in \mathcal{S} \cup \mathcal{T}}
   \Bigl(\alpha_{p,n} - \tilde{\alpha}_{p,n}\Bigr) \mathbf{v}_n
   \;+\; \sum_{n \in \mathcal{M}} \alpha_{p,n} \mathbf{v}_n \\
&= \sum_{n \in \mathcal{S} \cup \mathcal{T}}
   \frac{\tilde{Z}-Z}{Z \tilde{Z}}
   e^{\mathbf{q}_p^{\top} \mathbf{k}_n / \sqrt{d_k}} \, \mathbf{v}_n
   \;+\; \sum_{n \in \mathcal{M}} \alpha_{p,n} \mathbf{v}_n.
\end{align}
Since $Z - \tilde{Z} = \sum_{n \in \mathcal{M}}
e^{\mathbf{q}_p^{\top} \mathbf{k}_n / \sqrt{d_k}}$, the first term is a weighted
average over $\mathcal{S} \cup \mathcal{T}$ with total weight
$\frac{Z-\tilde{Z}}{Z}$, and the second term is a weighted average over
$\mathcal{M}$ with the same total weight.  Because the softmax weights sum to
$\frac{Z-\tilde{Z}}{Z}$ for the discarded mass, and $\|\mathbf{v}_n\|_2 \le 1$
after layer-normalization, the $\ell_2$ norm of the error is bounded by that
discarded mass, yielding Equation~\eqref{eq:rope-bound}.
\end{proof}

\textbf{Implication.}  Lemma~\ref{lem:rope} shows that accuracy degradation is
not caused by \textit{positional misalignment} but solely by the absence of the
middle tokens themselves.  By choosing a sufficiently large Tail window $L_t$,
the discarded mass becomes negligible for retrieval tasks that depend on
recent context, which explains the $100\%$ NIAH recall observed in our
experiments.

\subsection{Memory Bound Theorem}
\label{subsec:memory-bound}

\begin{theorem}[Constant Steady-State Memory]
\label{thm:memory}
For a model with $L$ layers, $h$ KV heads, and head dimension $d$,
the steady-state peak GPU memory of \texttt{Hetero-KVCache-Optimizer} is
\begin{equation}
M_{\text{steady}}(N) \;\le\;
M_{\text{weights}} \;+\;
2 \cdot L \cdot h \cdot (L_s + L_t) \cdot d \cdot \text{sizeof(dtype)}
\;+\; M_{\text{act}},
\label{eq:memory-bound}
\end{equation}
which is strictly $O(L_s + L_t)$ and completely independent of the input
sequence length $N$.  Here $L_s$ and $L_t$ are the Sink and Tail budgets, and
$M_{\text{act}}$ denotes activation scratch memory.
\end{theorem}

\begin{proof}
During the decode phase, the cache manager maintains exactly one physical
HBM buffer per layer.  By construction of the in-place rolling mechanism,
the buffer size is capped at $L_s + L_t$ tokens; when a new token arrives,
the oldest Tail token is overwritten in-place rather than re-allocated.
Each token stores a key and a value vector of shape $(h, d)$.
With $L$ layers, the total HBM-resident KV footprint is
$2 \cdot L \cdot h \cdot (L_s + L_t) \cdot d \cdot \text{sizeof(dtype)}$.

The transient prefill tensors are explicitly deleted and garbage-collected
after each chunk, so they do not contribute to steady-state consumption.
Overflow tokens are compressed to 4-bit and moved to DRAM (CPU pinned memory),
which does not count against GPU VRAM.  Therefore $M_{\text{steady}}(N)$
is constant in $N$ and bounded by Equation~\eqref{eq:memory-bound}.
\end{proof}

\textbf{Closed-form VRAM formula.}  For the concrete configuration used in our
evaluation ($L=28$, $h=4$, $d=128$, BF16 dtype, $L_s=64$, $L_t=8192$):
\begin{equation}
M_{\text{kv}} \;=\;
2 \times 28 \times 4 \times (64 + 8192) \times 128 \times 2\;\text{bytes}
\;\approx\; 0.454\;\text{GB},
\end{equation}
matching the empirical measurement reported in Table~\ref{tab:memory_comparison}.

\subsection{Bandwidth-Limited Latency Model}
\label{subsec:latency-model}

When the PCIe link between CPU DRAM and GPU HBM becomes the bottleneck,
the Time-Per-Output-Token (TPOT) is dominated by the latency of swapping in
the decompressed KV blocks needed for the next decode step.

\begin{theorem}[TPOT under PCIe Bandwidth Limit]
\label{thm:tpot}
Let $B$ be the effective PCIe bandwidth (GB/s) and $r$ the INT4 compression ratio
($r \approx 4.5$ for group size $128$).  During decode, each new token
may trigger the eviction of one Tail token to DRAM and the prefetch of another.
The per-token PCIe transfer size for layer $i$ is
\begin{equation}
\Delta_i \;=\;
\frac{2 \cdot h \cdot d \cdot \text{sizeof(BF16)}}{r}
\;=\;
\frac{4 \cdot h \cdot d}{r}\;\text{bytes}.
\label{eq:delta-i}
\end{equation}
With $L$ layers, the bandwidth-limited TPOT is bounded by
\begin{equation}
\text{TPOT}_{\text{PCIe}} \;\le\;
\frac{L \cdot \Delta_i}{B}
\;+\;
\text{TPOT}_{\text{compute}}.
\label{eq:tpot-bound}
\end{equation}
\end{theorem}

\begin{proof}
The rolling-window update evicts exactly one token per layer when the Tail
buffer is full.  That token is compressed to INT4 and offloaded to DRAM,
requiring $\Delta_i$ bytes of PCIe egress.  Symmetrically, if a prefetch is
scheduled for the next attention computation, an equal amount of data must
traverse the link in the opposite direction.  The worst-case PCIe time per
token is therefore the total bytes moved divided by the bidirectional bandwidth
$B$.  Adding the compute-bound TPOT (matrix multiplications, softmax, etc.)
gives the overall upper bound in Equation~\eqref{eq:tpot-bound}.
\end{proof}

\textbf{Predicted vs. measured values.}  Substituting our default configuration
($L=28$, $h=4$, $d=128$, $r=4.5$) into Equations~\eqref{eq:delta-i}
and~\eqref{eq:tpot-bound} yields:
\begin{equation}
\Delta_i
= \frac{4 \times 4 \times 128}{4.5}
\approx 455\;\text{bytes per layer}.
\end{equation}
On a full-bandwidth server link (PCIe Gen4 x16, $B \approx 32\;\text{GB/s}$):
\begin{equation}
\text{TPOT}_{\text{PCIe}}
\le \frac{28 \times 455}{32 \times 10^{9}} + 33\;\text{ms}
\approx 0.40\;\text{ms} + 33\;\text{ms}
\approx 33.4\;\text{ms}.
\end{equation}
On a constrained edge link (PCIe Gen4 x8, $B \approx 16\;\text{GB/s}$):
\begin{equation}
\text{TPOT}_{\text{PCIe}}
\le \frac{28 \times 455}{16 \times 10^{9}} + 33\;\text{ms}
\approx 0.80\;\text{ms} + 33\;\text{ms}
\approx 33.8\;\text{ms},
\end{equation}
and on a mobile edge link ($B \approx 8\;\text{GB/s}$):
\begin{equation}
\text{TPOT}_{\text{PCIe}}
\le \frac{28 \times 455}{8 \times 10^{9}} + 33\;\text{ms}
\approx 1.60\;\text{ms} + 33\;\text{ms}
\approx 34.6\;\text{ms}.
\end{equation}
These predictions confirm that the prefetch overhead remains below $5\%$ of
total TPOT even under severe bandwidth constraints, validating the viability
of the tiered storage abstraction for edge deployment.
